% TOP_PROBLEM: {"id":30007515,"problem_number":"LOCAL-30007515","title":"Sufficient heterogeneity in macro models","category_id":21,"category":{"id":21,"name":"economics","display_name":"Economics","description":"Open economic theory statements, published research agendas, and proposed scoped specifications; question provenance and historical priority are qualified per record.","slug":"economics","order_index":21,"created_at":"2026-10-08T00:00:00Z"},"status":"research_question","external_url":null,"legacy_ids":[],"published":false,"statement":"Find the smallest economically interpretable distributional state that approximates heterogeneous-agent policy responses uniformly over a restricted model and intervention class.","economics":{"original_id":4,"field":"Business cycles and macro dynamics","kind":"theoretical","formalization_basis":"adapted_research_question","source_status":"research_agenda","priority_status":"earliest_found","priority_note":"This is the earliest explicit relevant statement verified in this audit, not a claim of original priority; older literature and earlier versions may contain antecedents.","earliest_verified_reference":{"authors":["Janet L. Yellen"],"title":"Macroeconomic Research After the Crisis","year":2016,"url":"https://www.federalreserve.gov/newsevents/speech/yellen20161014a.htm","doi":null,"locator":"Section “Heterogeneity,” paragraphs on aggregate demand and financial constraints","evidence_summary":"The speech explicitly asks how individual differences affect aggregate demand and says heterogeneous financial constraints and firm linkages remain insufficiently understood empirically and theoretically."},"current_reference":{"authors":["Adrien Auclert","Matthew Rognlie","Ludwig Straub"],"title":"Fiscal and Monetary Policy with Heterogeneous Agents","year":2025,"url":"https://web.stanford.edu/~aauclert/annual_review_hank.pdf","doi":"10.1146/annurev-economics-091624-044646","locator":"Section 5, pp. 19–21 of the April 2025 author manuscript; “Nonlinearities” and “Optimal policy”","evidence_summary":"The review says empirical size and sign asymmetries remain debated. Its optimal-policy discussion identifies unresolved computation and equity–efficiency issues and warns that unrestricted monetary–fiscal Ramsey problems with standard preferences need not have a well-defined steady state."},"formalization_note":"A uniform approximation problem is newly specified to make the heterogeneity agenda precise; the sources do not publish this inequality as a conjecture."},"statusQualification":"Published question or research agenda verified at the cited locator. The mathematical specification is adapted for this catalogue and includes additional assumptions; source publication does not establish first-ever priority or certify this exact model remains unsolved. A uniform approximation problem is newly specified to make the heterogeneity agenda precise; the sources do not publish this inequality as a conjecture.","statusReviewedAt":"2026-10-08"}
% FIRST_POSTED: 2026-10-08
% ENTRY_KIND: statement_only
% !TeX program = lualatex
\documentclass[11pt]{article}
\usepackage{fontspec}
\usepackage[a4paper,margin=25mm]{geometry}
\usepackage{amsmath,amssymb,mathtools}
\usepackage{xurl}
\usepackage[hidelinks]{hyperref}
\newcommand{\sref}[2]{\hyperlink{source:#1:#2}{[S#2]}}
\setlength{\emergencystretch}{3em}
\begin{document}
% problemId:30007515
\section*{Sufficient heterogeneity in macro models}\label{30007515}
\subsection{Definitions and mathematical statement}
Let \(\mathcal M\) be a declared class of heterogeneous-agent economies sharing specified preferences, borrowing constraints, and production technologies, while their admissible asset, income, and firm-productivity distributions vary. For model \(M\in\mathcal M\), let \(\mathcal D_M\) be its admissible distribution set, \(\mu\) denote its cross-sectional distribution and \(R_M(\mu,a,h)\) its equilibrium response of output, consumption, employment, or inflation at horizon \(h\) to intervention \(a\). Interventions belong to a compact, explicitly bounded set \(\mathcal A\); horizons satisfy \(0\le h\le H\). Fix a response norm and accuracy \(\epsilon>0\).
Seek an economically interpretable statistic \(T(\mu)\in\mathbb R^d\), a transition law for that statistic, and a reduced equilibrium response \(\widehat R_M\) such that
\[
\sup_{M\in\mathcal M,\mu\in\mathcal D_M,a\in\mathcal A,\,0\le h\le H}
\|R_M(\mu,a,h)-\widehat R_M(T(\mu),a,h)\|\le\epsilon.
\]
Determine the smallest attainable dimension \(d\), or defensible upper and lower bounds, within a specified family of candidate statistics such as liquid-wealth bins, income exposure, and firm financing constraints. Explain which distributions and interventions invalidate a proposed compression. Establish approximation error separately from parameter-estimation error, and compare the compressed model with microeconomic evidence and held-out policy episodes. Without restrictions on \(\mathcal M\), finite-dimensional sufficiency need not exist; identifying those restrictions is part of this scoped target.

\par\textbf{Formulation and scope.} A uniform approximation problem is newly specified to make the heterogeneity agenda precise; the sources do not publish this inequality as a conjecture.

\par\textbf{Open-question provenance.} An explicit question or research agenda was verified at the source locator. This does not by itself certify that every later version remains unresolved \sref{30007515}{1}.

\par\textbf{Historical priority.} This is the earliest explicit relevant statement verified in this audit, not a claim of original priority; older literature and earlier versions may contain antecedents.

\subsection{Short English statement}
Find the smallest economically interpretable distributional state that approximates heterogeneous-agent policy responses uniformly over a restricted model and intervention class.

\subsection{Sources}
\begin{itemize}\raggedright
\item[\textnormal{[S1]}]\hypertarget{source:30007515:1}{} Janet L. Yellen. 2016. Macroeconomic Research After the Crisis. Section “Heterogeneity,” paragraphs on aggregate demand and financial constraints.
\url{https://www.federalreserve.gov/newsevents/speech/yellen20161014a.htm}.
{\small\textit{Earliest verified question/agenda reference; historical priority qualified below. The speech explicitly asks how individual differences affect aggregate demand and says heterogeneous financial constraints and firm linkages remain insufficiently understood empirically and theoretically.}}

\item[\textnormal{[S2]}]\hypertarget{source:30007515:2}{} Adrien Auclert, Matthew Rognlie, Ludwig Straub. 2025. Fiscal and Monetary Policy with Heterogeneous Agents. Section 5, pp. 19–21 of the April 2025 author manuscript; “Nonlinearities” and “Optimal policy”.
\url{https://web.stanford.edu/~aauclert/annual_review_hank.pdf}.
DOI: 10.1146/annurev-economics-091624-044646.
{\small\textit{Supporting or current literature; see evidence qualification. The review says empirical size and sign asymmetries remain debated. Its optimal-policy discussion identifies unresolved computation and equity–efficiency issues and warns that unrestricted monetary–fiscal Ramsey problems with standard preferences need not have a well-defined steady state.}}
\end{itemize}
\end{document}
